\begin{thema}{Β}
Έστω ένα ορθογώνιο παράθυρο με ένα ημικύκλιο ενσωματωμένο στην πάνω πλευρά του.
\begin{center}
\begin{tikzpicture}[scale=0.8]
    \draw[thick, fill=maincolor!10] (0, 0) -- (4, 0) -- (4, 3) arc (0:180:2) -- cycle;
    \draw[thick, dashed] (0, 3) -- (4, 3);
    \draw[<->, >=latex] (0, -0.3) -- (4, -0.3) node[midway, below] {$2x$};
    \draw[<->, >=latex] (-0.3, 0) -- (-0.3, 3) node[midway, left] {$y$};
    \draw[<->, >=latex] (4.3, 0) -- (4.3, 3) node[midway, right] {$y$};
    \filldraw (2, 3) circle (1.5pt) node[below] {};
    \draw[->, >=latex] (2, 3) -- (2.8, 4.83) node[midway, right] {$x$};
\end{tikzpicture}
\end{center}
Το συνολικό μήκος του πλαισίου του παραθύρου (δηλαδή οι τρεις ευθύγραμμες πλευρές του ορθογωνίου και το τόξο του ημικυκλίου, χωρίς την ενδιάμεση βάση) είναι $10$ μέτρα. Θεωρούμε το πλάτος του παραθύρου ίσο με $2x$ και το ύψος του ορθογώνιου μέρους ίσο με $y$, όπου $x, y > 0$.
\begin{erwthma}
  \item Να αποδείξετε ότι το εμβαδόν ολόκληρου του παραθύρου (το εμβαδόν του ορθογωνίου συν το εμβαδόν του ημικυκλίου) ως συνάρτηση του $x$ δίνεται από τον τύπο $E(x) = 10x - 2x^2 - \frac{\pi}{2}x^2$ για $x \in \left(0, \frac{10}{2+\pi}\right)$.
  \item Να μελετήσετε την $E(x)$ ως προς τη μονοτονία και να βρείτε την τιμή του $x$ για την οποία το εμβαδόν μεγιστοποιείται. Τι παρατηρείτε για τη σχέση του $x$ με το $y$ στη θέση αυτή?
  \item Να βρείτε το σύνολο τιμών της $E(x)$.
  \item Με βάση το σύνολο τιμών, να εξετάσετε πόσα διαφορετικά παράθυρα μπορούν να κατασκευαστούν που να έχουν συνολικό εμβαδόν $5 \; \tau.\mu.$.
\end{erwthma}
\end{thema}
